\section*{الفصل \ref{ch:g10:concentration-and-dilution} --- التركيز والتخفيف}

\begin{solution}{exo:g10:concentration-and-dilution:1}
$C_m = \qty{5.0}{g} / \qty{0.250}{L} = \qty{20}{g/L}$.
\end{solution}

\begin{solution}{exo:g10:concentration-and-dilution:2}
$c = \qty{0.050}{mol} / \qty{0.200}{L} = \qty{0.25}{mol/L}$.
\end{solution}

\begin{solution}{exo:g10:concentration-and-dilution:3}
$M(\ce{CuSO4}) = 63.5 + 32.1 + 4 \times 16.0 = \qty{159.6}{g/mol}$;
$n = 0.10 \times 0.1000 = \qty{0.0100}{mol}$; $m = 0.0100 \times 159.6 =
\qty{1.60}{g}$.
\end{solution}

\begin{solution}{exo:g10:concentration-and-dilution:4}
$F = 1.0 / 0.050 = 20$; $V_0 = \qty{100.0}{mL} / 20 = \qty{5.0}{mL}$.
\end{solution}

\begin{solution}{exo:g10:concentration-and-dilution:5}
$C_m = c \times M = 0.10 \times 58.5 = \qty{5.85}{g/L}$، أي نحو
\qty{5.9}{g/L}.
\end{solution}

\begin{solution}{exo:g10:concentration-and-dilution:6}
$V_0 = c_1 V_1 / c_0 = 0.020 \times 250.0 / 0.50 = \qty{10.0}{mL}$
(\omterm{def:g10:concentration-and-dilution:dilution}{معامل التخفيف} 25). ماصّة معيارية من \qty{10.0}{mL} وحوجلة معيارية من
\qty{250.0}{mL}.
\end{solution}

\begin{solution}{exo:g10:concentration-and-dilution:7}
للحوجلة المعيارية خط واحد، صُنع لحجم واحد وبدقة تبلغ جزءًا صغيرًا من الواحد في
المئة؛ أما تدريجات الكأس فليست إلا دليلًا تقريبيًا. وكذلك تصرف الماصّة
المعيارية حجمًا واحدًا بدقة كبيرة، بينما المخبار المدرّج أعرض ويُقرأ بدقة
أقل.
\end{solution}

\begin{solution}{exo:g10:concentration-and-dilution:8}
بين \qty{0.02}{mol/L} و\qty{0.04}{mol/L}: أي نحو \qty{0.03}{mol/L}.
ولدقة أكبر، أضف محاليل معيارية بين هذين، مثلًا عند
0.025 و0.030 و\qty{0.035}{mol/L}: أي سلّمًا أدق حيث تلزم الدقة.
\end{solution}

\begin{solution}{exo:g10:concentration-and-dilution:9}
$M(\ce{C6H12O6}) = \qty{180.0}{g/mol}$; $c = 50 / 180.0 =
\qty{0.28}{mol/L}$.
\end{solution}

\begin{solution}{exo:g10:concentration-and-dilution:10}
$M = \qty{342.0}{g/mol}$، إذن $n = \qty{1.00}{mol}$ و$c =
\qty{1.00}{mol/L}$. وبعد \omterm{def:g10:concentration-and-dilution:dilution}{التخفيف} عشر مرات: $c = \qty{0.100}{mol/L}$، أي
$C_m = 0.100 \times 342.0 = \qty{34.2}{g/L}$. وملعقة من
\qty{0.020}{L} تحوي $34.2 \times 0.020 = \qty{0.68}{g}$ من السكروز.
\end{solution}

\begin{solution}{exo:g10:concentration-and-dilution:11}
(أ) الكتلة نفسها من \omterm{def:g6:solutions-and-solubility:solute}{المذاب} في ضعف الحجم: \qty{6.0}{g/L}.
(ب) ينقص الحجم إلى النصف تقريبًا، وكتلة \omterm{def:g6:solutions-and-solubility:solute}{المذاب} لا تتغير: نحو
\qty{24}{g/L}.
\end{solution}

\begin{solution}{exo:g10:concentration-and-dilution:12}
\begin{enumerate}
  \item ثلاثة تخفيفات بعشرة أضعاف تساوي تخفيفًا بمعامل $10^3$:
        $\qty{0.10}{mol/L} / 1000 = \qty{1.0e-4}{mol/L}$.
  \item \omterm{def:g10:concentration-and-dilution:dilution}{التخفيف} الواحد بألف ضعف يتطلب إما حجمًا ضئيلًا من \omterm{def:g10:concentration-and-dilution:dilution}{المحلول الأم}
        (\qty{0.1}{mL} في \qty{100}{mL})، يتعذر قياسه بدقة بماصّة، وإما
        حوجلة كبيرة جدًا (\qty{1}{mL} في \qty{1}{L}). أما ثلاثة تخفيفات
        بعشرة أضعاف فتستعمل ماصّات وحوجلات
        عادية.
\end{enumerate}
\end{solution}

\begin{solution}{exo:g10:concentration-and-dilution:13}
\begin{enumerate}
  \item $M(\ce{CaCl2}) = 40.1 + 2 \times 35.5 = \qty{111.1}{g/mol}$؛
        $n = 11.1 / 111.1 = \qty{0.100}{mol}$؛ $c = 0.100 / 0.5000 =
        \qty{0.200}{mol/L}$.
  \item كل \ce{CaCl2} يعطي \omterm{def:g9:ions:ion}{أيون} \ce{Ca^2+} واحدًا وأيوني \ce{Cl-}:
        $[\ce{Ca^2+}] = \qty{0.200}{mol/L}$ و$[\ce{Cl-}] =
        \qty{0.400}{mol/L}$.
\end{enumerate}
\end{solution}

\begin{solution}{exo:g10:concentration-and-dilution:14}
\begin{enumerate}
  \item $r = \qty{0.60}{cm}$، $h = \qty{0.20}{cm}$:
        $V = \pi \times 0.60^2 \times 0.20 = \qty{0.23}{cm^3}$، أي
        \qty{0.23}{mL} من الماء الزائد.
  \item \omterm{def:g6:solutions-and-solubility:solute}{المذاب} نفسه في \qty{100.23}{mL} بدلًا من \qty{100.00}{mL}:
        فالتركيز أقل من اللازم بنسبة $0.23 / 100.23 \approx
        \qty{0.23}{\percent}$.
\end{enumerate}
\end{solution}

\begin{solution}{exo:g10:concentration-and-dilution:15}
$n = 0.20 \times 0.100 + 0.10 \times 0.300 = 0.020 + 0.030 =
\qty{0.050}{mol}$ في \qty{0.400}{L}: $c = \qty{0.125}{mol/L}$. وليس هو
المتوسط البسيط، \qty{0.15}{mol/L}: \omterm{def:g2:mixing-and-dissolving:dissolve}{فالمحلول} المخفف أكثر بثلاث مرات، فيكون
\omterm{def:g6:pure-substances-and-mixtures:mixture}{الخليط} أقرب إلى \qty{0.10}{mol/L}. إنه المتوسط المرجّح
بالحجوم.
\end{solution}

\begin{solution}{pb:g10:concentration-and-dilution:1}
\textbf{1.} \qty{0.9}{g} في \qty{0.100}{L}: $C_m = \qty{9.0}{g/L}$.

\textbf{2.} \qty{9.0}{g} في \qty{1000}{mL}، أي \qty{9.0}{mg} لكل
ميليلتر.

\textbf{3.} $m = 9.0 \times 0.500 = \qty{4.5}{g}$.

\textbf{4.} نعم: نحو \qty{9}{g} في اللتر مقابل نحو \qty{360}{g} في
الكيلوغرام من الماء عند الإشباع، أي أقل بنحو أربعين مرة.

\textbf{5.} $9.0 \times 1.0 = \qty{9.0}{g}$ من الملح.

\textbf{6.} $M(\ce{NaCl}) = 23.0 + 35.5 = \qty{58.5}{g/mol}$.

\textbf{7.} $c = 9.0 / 58.5 = \qty{0.154}{mol/L}$.

\textbf{8.} $n = 0.154 \times 0.500 = \qty{0.0769}{mol}$ (أو
$4.5 / 58.5$).

\textbf{9.} كل \ce{NaCl} يعطي \ce{Na+} واحدًا و\ce{Cl-} واحدًا: فكلاهما عند
\qty{0.154}{mol/L}.

\textbf{10.} $0.0769 \times \num{6.02e23} = \num{4.63e22}$ \omterm{def:g9:ions:ion}{أيون} صوديوم.

\textbf{11.} $F = 200 / 9.0 = 22.2$.

\textbf{12.} الغرامات \qty{4.5}{g} في الكيس تأتي كلها من \omterm{def:g2:mixing-and-dissolving:dissolve}{المحلول} المركّز:
$V_0 = \qty{4.5}{g} / \qty{200}{g/L} = \qty{0.0225}{L} =
\qty{22.5}{mL}$.

\textbf{13.} $V_0 = V_1 / F = 500 / 22.2 = \qty{22.5}{mL}$.

\textbf{14.} $c_0 = 200 / 58.5 = \qty{3.42}{mol/L}$؛
و$c_0 V_0 = 3.42 \times 0.0225 = \qty{0.0769}{mol}$
و$c_1 V_1 = 0.154 \times 0.500 = \qty{0.0769}{mol}$: متساويان.

\textbf{15.} نحو $500 - 22.5 = \qty{478}{mL}$ من الماء المعقم.

\textbf{16.} ماصّة مدرّجة من \qty{25}{mL} (أو سحاحة)، تُقرأ إلى
\qty{0.1}{mL}.

\textbf{17.} في حوجلة معيارية من \qty{500.0}{mL}، يُكمَّل فيها الحجم حتى الخط.

\textbf{18.} $C_m = 200 \times 23.0 / 500 = \qty{9.20}{g/L}$، أي أكثر تركيزًا
مما ينبغي بنسبة $0.20 / 9.0 \approx \qty{2.2}{\percent}$.

\textbf{19.} \qty{22.5}{mL} من \omterm{def:g2:mixing-and-dissolving:dissolve}{المحلول} المركّز («20 \%») لكل كيس من
\qty{500}{mL}.
\end{solution}
